\chapter{Sélection et lecture du modèle sur 2023}
\label{chap:validation-ml}

Le notebook 06 utilise 2023 pour choisir la régularisation, comparer la logistique à
la référence et examiner ses comportements. Les résultats présentés ici sont ceux
exportés dans le projet fourni, avec leurs métadonnées.

\section{Un réglage limité à quatre valeurs de C}
La grille est fixée à $C\in\{0{,}01;\,0{,}1;\,1;\,10\}$. La règle retient la plus
petite log-loss de validation et, en cas d'égalité exacte, le plus petit $C$.
La logistique n'est conservée que si elle fait strictement mieux que la référence.
\begin{table}[htbp]
\centering\small
\renewcommand{\arraystretch}{1.18}
\caption{Comparaison des quatre régularisations sur 2023 ; ligne retenue en gras.}\label{tab:ml-selection}
\begin{tabular*}{\textwidth}{@{\extracolsep{\fill}}lrrrr@{}}
\toprule
C & Log-loss & Brier & AP & ROC-AUC \\
\midrule
0,01 & 0,401516 & 0,127054 & 0,408526 & 0,768392 \\
0,1 & 0,400823 & 0,127067 & 0,408773 & 0,769805 \\
\textbf{1} & \textbf{0,400797} & \textbf{0,127102} & \textbf{0,408514} & \textbf{0,769915} \\
10 & 0,400805 & 0,127110 & 0,408432 & 0,769911 \\
\bottomrule
\end{tabular*}
\sourcecalcul{notebook 06 ; même représentation et même apprentissage pour les quatre valeurs de C}
\end{table}


La valeur retenue est \textbf{$C=1$}. Les résultats sont néanmoins très proches pour
$C=0{,}1$, $1$ et $10$. Le protocole départage ces candidats, mais ces écarts ne
justifient pas de présenter $C=1$ comme une valeur universellement optimale.
La grille ne devient pas plus fine après consultation du test.

\section{L'apport par rapport à la fréquence constante}
\begin{table}[htbp]
\centering\small
\renewcommand{\arraystretch}{1.18}
\caption{Scores de validation de la référence et du modèle retenu, 2023.}\label{tab:ml-validation}
\begin{tabular*}{\textwidth}{@{\extracolsep{\fill}}lrrrr@{}}
\toprule
Modèle & Log-loss & Brier & AP & ROC-AUC \\
\midrule
Référence constante & 0,4709 & 0,1474 & 0,1797 & 0,5000 \\
Régression logistique & 0,4008 & 0,1271 & 0,4085 & 0,7699 \\
\bottomrule
\end{tabular*}
\sourcecalcul{notebook 06 ; comparaison après sélection sur cette même année}
\end{table}

La log-loss passe de 0,4709 à 0,4008, soit une réduction relative de 14,88\,\%.
Le Brier diminue de 13,76\,\%. L'AP atteint 0,4085 et la ROC-AUC 0,7699.
Le modèle apporte donc une information utile au sens des critères retenus~: les
caractéristiques distinguent mieux les issues que la seule fréquence historique.

Ce résultat ne signifie ni que 77\,\% des usagers seraient correctement classés,
ni que l'on disposerait d'un outil opérationnel. La ROC-AUC décrit un classement~;
la log-loss et le Brier évaluent les probabilités sur cette population.
Le passage à une décision nécessite encore un seuil et une justification de ses conséquences.

\section{Calibration et compromis précision-rappel}
La calibration compare des probabilités annoncées à des proportions observées pour
des groupes de prédictions~\cite{sklearn-calibration}. Le notebook utilise dix
intervalles fixes, de largeur 0,1. Une courbe proche de la diagonale est compatible
avec un bon accord dans ces intervalles~; elle ne prouve pas une calibration exacte
pour chaque profil.
\figurerapportML[0.98]{06_01_calibration_validation.png}{Calibration descriptive et distribution des probabilités sur 2023.}{fig:calibration-val}{06}

L'accord est imparfait. Les probabilités élevées sont notamment surestimées~: dans
l'intervalle $]0{,}7;0{,}8]$, la moyenne prédite est de 74,06\,\% contre 60,43\,\%
observés, pour 988 usagers. Les derniers intervalles reposent sur moins de cas et
demandent une lecture prudente. Aucune recalibration n'est sélectionnée dans cette version.

\figurerapportML[0.88]{06_02_precision_rappel_validation.png}{Courbe précision-rappel sur 2023 et proportion positive de référence.}{fig:pr-val}{06}
La précision mesure la part des cas positifs parmi les usagers classés positifs.
Le rappel mesure la part des cas positifs retrouvés. La courbe montre que couvrir une
plus grande fraction des décès ou hospitalisations s'accompagne ici d'une précision
plus faible. Le seuil 0,5 reste un exemple fixé à l'avance~; il n'est pas choisi
pour maximiser une métrique sur 2023.

\section{Les résultats moyens masquent des différences}
La proportion positive de 2023 est de 16,79\,\% chez les conducteurs, 16,86\,\%
chez les passagers et 32,03\,\% chez les piétons. Le modèle retrouve assez bien ces
niveaux moyens, mais son classement est moins performant au sein du groupe des piétons~:
la ROC-AUC est de 0,6074, contre 0,7919 chez les conducteurs.
\figurerapportML[0.88]{06_03_diagnostics_categories_validation.png}{Proportions observées et probabilités moyennes par catégorie d'usager en 2023.}{fig:groupes-val}{06}

Un bon accord sur la moyenne d'un groupe ne suffit donc pas à distinguer correctement
ses individus. Les effectifs, les cas positifs et les scores par groupe complètent
nécessairement le graphique des moyennes.

Un autre diagnostic concerne les informations manquantes. En 2023, les 2\,359 usagers
dont l'âge ou le sexe n'est pas renseigné ne comptent aucun cas positif. Le modèle
leur attribue de très faibles probabilités. Cette structure peut traduire un mécanisme
de renseignement plutôt qu'une relation transportable entre absence d'information
et gravité. Sur les seuls usagers dont ces deux champs sont renseignés, sans
réentraînement, la log-loss reste meilleure que la référence~: 0,4089 contre 0,4765.
Le gain ne repose donc pas uniquement sur ce groupe, sans que cette restriction
corrige les biais de collecte.

\section{Relier les coefficients aux comparaisons statistiques}
L'encodage conserve toutes les modalités. La comparaison de deux modalités A et R
d'un même champ se lit par la différence de leurs coefficients~:
\[
  \mathrm{OR}_{A/R}=\exp(\widehat\beta_A-\widehat\beta_R).
\]
Ce rapport de cotes est conditionnel aux autres caractéristiques du modèle et dépend
de sa représentation et de la pénalisation. Il ne s'agit pas d'un rapport de proportions.
\begin{table}[htbp]
\centering\small
\renewcommand{\arraystretch}{1.18}
\caption{Contrastes de coefficients de la logistique retenue.}\label{tab:ml-contrastes}
\begin{tabularx}{\textwidth}{@{}Lrr@{}}
\toprule
Comparaison A / R & Écart de log-cotes & Rapport de cotes \\
\midrule
75 ans ou plus / 35-44 ans & 1,369 & 3,932 \\
Piéton / conducteur & 1,608 & 4,995 \\
Nuit sans éclairage / plein jour & 0,219 & 1,245 \\
Hors / en agglomération & 1,203 & 3,329 \\
\bottomrule
\end{tabularx}
\sourcecalcul{notebook 06 ; rapports pénalisés, conditionnels au modèle, sans intervalle ni interprétation causale}
\end{table}


Pour les piétons, le contraste vaut environ 4,99 en cotes, à caractéristiques du modèle
maintenues identiques. Il ne signifie pas «~cinq fois la probabilité~» et ne mesure
pas l'effet causal d'être piéton. Le rapport brut de proportions de la première partie
répondait à une autre question. La différence entre ces lectures est particulièrement
visible pour la lumière~: le contraste de nuit sans éclairage devient beaucoup plus
modéré une fois les autres variables incluses. Cela ne permet pas d'attribuer l'écart
à une cause précise.

Aucun intervalle ni test de significativité des coefficients n'est fourni. Cette
lecture reste exploratoire, avec les limites du modèle additif et des profils parfois
peu représentés dans les données.

\FloatBarrier
\bilan{La logistique est retenue sur 2023, avec $C=1$. Elle améliore les scores de
probabilité et de classement, mais les diagnostics mettent déjà en évidence une
surestimation des fortes probabilités, un classement plus faible chez les piétons et
un signal lié aux informations absentes. La configuration est figée pour le test.}
